%============================================================================
% LICENSING NOTICE
%============================================================================
% This WIP file, upon integration into guide-v.tex, becomes part of the
% TMS/DMS/CMS Usage Guide and is released under CC BY-NC 4.0.
% For details, see LICENSE in the repository root.
%============================================================================

%============================================================================
% FALCON BMS TMS/DMS/CMS HOTAS GUIDE
% WIP FILE TEMPLATE — FINAL (Briefing)
%============================================================================
% IMPORTANTE: Este é um TEMPLATE padronizado para arquivos WIP (Work-In-Progress)
% que serão integrados ao guide.tex.
% Nomenclatura: Siga wip-naming
% Padrão: section-C{N}-S{M}[-S{K}]-{titulo}-{status}-{data}.tex
% Exemplo: section-C5-S2-cms-actuation-hotas-tables-final-2026-01-10.tex
% Status: dev | review | final | alpha | approved | deprecated
% Note: 'alpha' is the parallel pre-release flow (optional) — see wip-naming.md §0.5
% Locação: WIP/ (ativo) | ARCHIVE/ (aprovado/descartado)

%============================================================================
% PREAMBLE COMPLETO — TMS/DMS/CMS Usage Guide for Falcon BMS 4.38.1
% Gerado: 25 March 2026
%============================================================================

\documentclass[11pt, a4paper, twoside]{report}

% --------------------------------------------------------------------------
% BASIC ENCODING AND LANGUAGE
% --------------------------------------------------------------------------

\usepackage[utf8]{inputenc}
\usepackage[T1]{fontenc}
\usepackage[english]{babel}

% --------------------------------------------------------------------------
% FONTS AND MICROTYPOGRAPHY
% --------------------------------------------------------------------------

\usepackage{lmodern}
\usepackage{microtype}

% --------------------------------------------------------------------------
% PAGE GEOMETRY AND LAYOUT
% --------------------------------------------------------------------------

\usepackage{geometry}
\geometry{a4paper, left=2.0cm, right=2.0cm, top=2.5cm, bottom=2.5cm}
\usepackage{setspace}
\onehalfspacing

% --------------------------------------------------------------------------
% VERSION CONTROL MACROS
% --------------------------------------------------------------------------

\newcommand{\docversion}{0.4.2.0.alpha.3}
\newcommand{\docbuild}{20260325}
\newcommand{\docstartdate}{05 January 2026}
\newcommand{\docenddate}{20260325}
\newcommand{\chapterscompletedof}{1/6}
\newcommand{\tablesfilledpct}{N/A}
\newcommand{\fulldocversion}{\docversion+\docbuild}

% --------------------------------------------------------------------------
% COLORS AND LINKS AND PDF SETTINGS
% --------------------------------------------------------------------------

\usepackage[table]{xcolor}
\definecolor{linkblue}{HTML}{004488}
\definecolor{linkred}{HTML}{882222}
\definecolor{headerblue}{HTML}{003366}
\definecolor{rowgray}{HTML}{F5F5F5}
\definecolor{subheadgray}{HTML}{E0E0E0}
\usepackage{soul}
\usepackage[pdfencoding=auto, psdextra, colorlinks=true, linkcolor=linkblue, citecolor=linkred, urlcolor=linkblue, breaklinks=true]{hyperref}
\usepackage{hyperxmp}
\hypersetup{
	pdftitle={TMS, DMS and CMS Usage Guide for Falcon BMS},
	pdfauthor={Carlos "Metal" Nader},
	pdfsubject={Flight Simulation - Falcon BMS HOTAS Reference},
	pdfkeywords={Falcon BMS, F-16, HOTAS, TMS, DMS, CMS, Flight Simulation},
	pdfcreator={pdfLaTeX (TeXstudio)},
	pdfproducer={TeXstudio},
	pdflang={en-US},
	pdfcopyright={CC BY-NC 4.0 · Carlos "Metal" Nader},
	pdflicenseurl={https://creativecommons.org/licenses/by-nc/4.0/},
	pdfcontacturl={https://github.com/carlos-nader/tms-dms-cms-usage-guide},
	pdfdisplaydoctitle=true,
	bookmarksnumbered=true,
	bookmarksopen=true,
	bookmarksdepth=1,
	pdfstartview=FitH,
}
\pdfinfo{/CreationDate (D:\docbuild) /ModDate (D:\docbuild)}
\usepackage{bookmark}

% --------------------------------------------------------------------------
% CAPTION SETUP BY TYPE
% --------------------------------------------------------------------------

\usepackage{caption}

% GLOBAL PATTERN
\captionsetup{font=small, labelfont=bf, justification=centering, singlelinecheck=true}
% FIGURES
\captionsetup[figure]{font=footnotesize, labelfont=bf, justification=centering, singlelinecheck=true}
% % TABLE/LONGTABLE/hotastable:
\captionsetup[table]{font=small, labelfont=bf, justification=centering, singlelinecheck=true}

% --------------------------------------------------------------------------
% HEADERS AND FOOTERS (IMPROVED for report + twoside)
% --------------------------------------------------------------------------

\usepackage{fancyhdr}
\setlength{\headheight}{25pt}
\pagestyle{fancy}
\fancyhf{}
\fancyhead[LO,RE]{\small\textit{\leftmark}}
\fancyhead[RO,LE]{\small\thepage}
\fancyfoot{}
\renewcommand{\headrulewidth}{0.4pt}
\renewcommand{\footrulewidth}{0pt}

% --------------------------------------------------------------------------
% CHAPTER FORMATTING AND SPACING (via titlesec)
% --------------------------------------------------------------------------

\usepackage{titlesec}

% --------------------------------------------------------------------------
% CHAPTER - Nível 0 (Destaque Máximo)
% --------------------------------------------------------------------------
% Shape: display → standalone em nova linha
% Font: \Large\bfseries --- label e título ambos em bold para destaque máximo
% Numeração: SIM
% Espaçamento: 15pt antes, 28pt depois
% --------------------------------------------------------------------------

\titleformat{\chapter}[display]
{\normalfont\Large\bfseries}           % Formato: LARGE + bold (label)
{\chaptertitlename~\thechapter}        % Rótulo: "Chapter 1" (bold também)
{20pt}                                 % Espaço vertical antes do corpo
{\Large\bfseries}                      % Corpo do título: LARGE + bold

\titlespacing{\chapter}
{0pt}       % Margem esquerda (sem indent)
{15pt}      % Espaço ANTES do título do capítulo
{28pt}      % Espaço DEPOIS do título do capítulo (AUMENTADO)
[0pt]       % Margem direita

% --------------------------------------------------------------------------
% SECTION - Nível 1
% --------------------------------------------------------------------------
% Shape: hang → label à esquerda, corpo indentado (compacto)
% Font: \large\bfseries (mesmo destaque que chapter para coerência)
% Numeração: SIM
% Espaçamento: 15pt antes, 8pt depois (REDUZIDO de 15pt/10pt anteriores)
% --------------------------------------------------------------------------

\titleformat{\section}[hang]
{\normalfont\large\bfseries}           % Formato: large + bold
{\thesection}                          % Rótulo: "1", "2", etc
{1em}                                  % Espaço entre rótulo e corpo
{}                                     % Corpo do título (herda do format)

\titlespacing{\section}
{0pt}       % Margem esquerda (sem indent)
{30pt}      % Espaço ANTES do título da seção
{8pt}       % Espaço DEPOIS do título da seção
[0pt]       % Margem direita

% --------------------------------------------------------------------------
% SUBSECTION - Nível 2
% --------------------------------------------------------------------------
% Shape: hang → label esquerda, corpo indentado
% Font: \normalsize\bfseries (redução visual, ainda robusto)
% Numeração: SIM
% Espaçamento: 12pt antes, 6pt depois (compacto)
% --------------------------------------------------------------------------

\titleformat{\subsection}[hang]
{\normalfont\normalsize\bfseries}      % Formato: tamanho normal + bold
{\thesubsection}                       % Rótulo: "1.1", "1.2", etc
{1em}                                  % Espaço entre rótulo e corpo
{}                                     % Corpo do título

\titlespacing{\subsection}
{0pt}       % Margem esquerda (sem indent)
{24pt}      % Espaço ANTES do título da subseção
{6pt}       % Espaço DEPOIS do título da subseção
[0pt]       % Margem direita

% --------------------------------------------------------------------------
% SUBSUBSECTION - Nível 3
% --------------------------------------------------------------------------
% Shape: hang → label esquerda, máximo recuo (bem compacto)
% Font: \normalsize\bfseries (mesmo tamanho que subsection, padrão)
% Numeração: SIM (configurable com secnumdepth)
% Espaçamento: 10pt antes, 4pt depois (bem comprimido)
% --------------------------------------------------------------------------

\titleformat{\subsubsection}[hang]
{\normalfont\normalsize\bfseries}      % Formato: tamanho normal + bold
{\thesubsubsection}                    % Rótulo: "1.1.1", "1.1.2", etc
{1em}                                  % Espaço entre rótulo e corpo
{}                                     % Corpo do título

\titlespacing{\subsubsection}
{0pt}       % Margem esquerda (sem indent)
{18pt}       % Espaço ANTES do título
{4pt}       % Espaço DEPOIS do título (bem comprimido)
[0pt]       % Margem direita

% --------------------------------------------------------------------------
% PARAGRAPH - Nível 4 (CRÍTICO: Ilusão Óptica Resolvida)
% --------------------------------------------------------------------------
% Shape: runin → inline, integrado ao texto (seu requisito)
% Font: \small\bfseries (SOLUÇÃO PARA ILUSÃO ÓPTICA DO NEGRITO)%
% Numeração: NÃO (padrão para \paragraph)
% Espaçamento: runin (sem espaço vertical antes, integrado ao texto)
% --------------------------------------------------------------------------

\titleformat{\paragraph}[runin]
{\normalfont\small\bfseries}           % Formato: SMALL + bold (compensa ilusão)
{}                                     % Rótulo: vazio (sem numeração)
{0em}                                  % Espaço entre rótulo e corpo (zero, runin)
{}                                     % Corpo do título

\titlespacing{\paragraph}
{0pt}       % Margem esquerda (sem indent)
{8pt}       % Espaço ANTES do título do parágrafo (ADICIONADO - cria separação)
{1em}       % Espaço DEPOIS do título (ADICIONADO - espaço após ":")
[0pt]       % Margem direita

% --------------------------------------------------------------------------
% SUBPARAGRAPH - Nível 5 (Preparação Futura)
% --------------------------------------------------------------------------
% Shape: runin → inline, integrado ao texto
% Font: \small (MENOS destaque que \paragraph, sem bold)
% Numeração: NÃO (padrão para \subparagraph)
% Espaçamento: runin (integrado)
% Uso: Se usado, aparecerá com destaque menor que \paragraph
% --------------------------------------------------------------------------

\titleformat{\subparagraph}[runin]
{\normalfont\small}                    % Formato: small, SEM bold (menos destaque)
{}                                     % Rótulo: vazio (sem numeração)
{0em}                                  % Espaço entre rótulo e corpo (zero, runin)
{}                                     % Corpo do título

\titlespacing{\subparagraph}
{0pt}       % Margem esquerda (sem indent)
{8pt}       % Espaço ANTES do título (ADICIONADO)
{1em}       % Espaço DEPOIS do título (ADICIONADO)
[0pt]       % Margem direita

% --------------------------------------------------------------------------
% TABLES AND MACROS
% --------------------------------------------------------------------------

\usepackage{booktabs}
\usepackage{array}
\usepackage{longtable}
\usepackage{tabularx}

% Custom Columns
\newcolumntype{L}[1]{>{\raggedright\arraybackslash}p{#1}}
\newcolumntype{C}[1]{>{\centering\arraybackslash}p{#1}}
\newcolumntype{R}[1]{>{\raggedleft\arraybackslash}p{#1}}

% Macro for Visual Reference Links
\newcommand{\imglink}[1]{\hspace{2pt}\hyperref[#1]{\scriptsize\textbf{[Fig]}}}

% --------------------------------------------------------------------------
% HOTAS LOT FILE HANDLE
% --------------------------------------------------------------------------
\makeatletter
\newcommand{\listofhotastables}{%
	\section*{List of HOTAS Tables}%
	\addcontentsline{toc}{section}{List of HOTAS Tables}%
	\@starttoc{hotas}%
}
\makeatother

% --------------------------------------------------------------------------
% HOTAS table environment
% Version: 2.1 (2026-02-10)
% --------------------------------------------------------------------------

\newenvironment{hotastable}[1]{%
	\footnotesize
	\setlength{\tabcolsep}{3pt}
	\renewcommand{\arraystretch}{1.35}
	\begin{longtable}{L{1.00cm} L{0.90cm} L{0.90cm} L{3.30cm} L{6.40cm} L{1.40cm} L{1.60cm}}
		\caption{#1}\label{table.\thetable}\\
		\noalign{\addtocontents{hotas}{\protect\contentsline{table}{\protect\numberline{\thetable}#1}{\thepage}{table.\thetable}}}%
		\rowcolor{headerblue}
		\multicolumn{1}{>{\centering\arraybackslash}p{1.00cm}}{\textbf{\color{white}Mode}} &
		\multicolumn{1}{>{\centering\arraybackslash}p{0.90cm}}{\textbf{\color{white}Dir.}} &
		\multicolumn{1}{>{\centering\arraybackslash}p{0.90cm}}{\textbf{\color{white}Act.}} &
		\multicolumn{1}{>{\centering\arraybackslash}p{3.30cm}}{\textbf{\color{white}Function}} &
		\multicolumn{1}{>{\centering\arraybackslash}p{6.40cm}}{\textbf{\color{white}Effect / Nuance}} &
		\multicolumn{1}{>{\centering\arraybackslash}p{1.40cm}}{\textbf{\color{white}Dash34}} &
		\multicolumn{1}{>{\centering\arraybackslash}p{1.60cm}}{\textbf{\color{white}Train.}} \\
		\endfirsthead
		\rowcolor{headerblue}
		\multicolumn{1}{>{\centering\arraybackslash}p{1.00cm}}{\textbf{\color{white}Mode}} &
		\multicolumn{1}{>{\centering\arraybackslash}p{0.90cm}}{\textbf{\color{white}Dir.}} &
		\multicolumn{1}{>{\centering\arraybackslash}p{0.90cm}}{\textbf{\color{white}Act.}} &
		\multicolumn{1}{>{\centering\arraybackslash}p{3.30cm}}{\textbf{\color{white}Function}} &
		\multicolumn{1}{>{\centering\arraybackslash}p{6.40cm}}{\textbf{\color{white}Effect / Nuance}} &
		\multicolumn{1}{>{\centering\arraybackslash}p{1.40cm}}{\textbf{\color{white}Dash34}} &
		\multicolumn{1}{>{\centering\arraybackslash}p{1.60cm}}{\textbf{\color{white}Train.}} \\
		\endhead
		\multicolumn{7}{r}{\small\emph{Continued on next page}}\\
		\endfoot
		\endlastfoot
	}{%
	\end{longtable}
}

% --------------------------------------------------------------------------
% SIMPLE REFERENCE MACROS FOR BMS DOCS
% --------------------------------------------------------------------------

\newcommand{\dashref}[1]{Dash-34~\S~#1}
\newcommand{\dashone}[1]{Dash-1~\S~#1}
\newcommand{\trnref}[1]{TRN~#1}
\newcommand{\trnman}{BMS Training Manual 4.38.1}
\newcommand{\bmsver}{Falcon BMS~4.38.1}
\newcommand{\dashrefs}[1]{\textit{TO 1F-16CMAM-34-1-1}, Dash-34, sections \texttt{#1}}

% --------------------------------------------------------------------------
% CROSS-REFERENCE MACROS (INTERNAL GUIDE)
% --------------------------------------------------------------------------

\newcommand{\secref}[1]{\hyperref[#1]{Section~\ref*{#1}}}
\newcommand{\chapref}[1]{\hyperref[#1]{Chapter~\ref*{#1}}}
\newcommand{\tabref}[1]{\hyperref[#1]{Table~\ref*{#1}}}
\newcommand{\figref}[1]{\hyperref[#1]{Figure~\ref*{#1}}}

% --------------------------------------------------------------------------
% GRAPHICS
% --------------------------------------------------------------------------

\usepackage{graphicx}
\graphicspath{{fig/}}
\usepackage{float}
\usepackage{enumitem}

% --------------------------------------------------------------------------
% TITLE
% --------------------------------------------------------------------------

\title{TMS, DMS and CMS Usage Guide for \bmsver}
\author{Carlos ``Metal'' Nader}
\date{Version \fulldocversion{} | Progress: Chapters \chapterscompletedof{} | Tables \tablesfilledpct{} | January 2026}

% ============================================================================
% DOCUMENT BEGIN
% ============================================================================

\begin{document}

	\maketitle

	\pagenumbering{roman}

	% --------------------------------------------------------------------------
	% TOC DEPTH CONFIGURATION
	% --------------------------------------------------------------------------
	\setcounter{tocdepth}{3}       % Show up to \subsubsection in TOC
	\setcounter{secnumdepth}{3}    % Number up to \subsubsection

	\newpage
	\tableofcontents
	\newpage

	% --------------------------------------------------------------------------
	% LIST OF HOTAS TABLES
	% --------------------------------------------------------------------------
	\phantomsection
	\listofhotastables
	\newpage

	\pagenumbering{arabic}

%============================================================================
% WIP FILE METADATA (NOT RENDERED IN PDF)
%============================================================================
% File Name: chapter-C2-style-rev-review-2026-03-27.tex
% Issue: #47 #44
% WIP Naming Convention: 2026-03-24
% Target Chapter: C2 (HOTAS Fundamentals)
% WIP Status: review
% Created: 2026-03-20
% Last Modified: 2026-03-27
% Integration Status: NOT YET INTEGRATED | INTEGRATION TARGET: alpha.3
%
% Narrative Completion: ~100% (all sections complete; pending final author review)
% Table Fill Status: N/A
%
% Notes:
% - Style revision pass for C2 — milestone v0.4.2.0, parent #44
% - #47 Task 1 (CMS orientation): WILL NOT IMPLEMENT — content covered by
%   existing §2.3.3 and C5 alpha.2
% - #47 Task 2 (TGP paragraph into §2.1.3): WILL NOT IMPLEMENT — TGP already
%   present as example in §2.1.3; detail belongs in C4 (TMS)
% - docs-web/c2-hotas-fundamentals.md updated to mirror WIP content (2026-03-27)

%============================================================================
% CONTENT BEGINS HERE (use \chapter{}, \section{}, \subsection{}, etc.)
%============================================================================

%============================================================================
% CHAPTER 2 — HOTAS FUNDAMENTALS
%============================================================================
\chapter{HOTAS Fundamentals}
\label{chap:C2}
The F-16 HOTAS (Hands On Throttle And Stick) controls consist of switches on the throttle and side-stick controller. The Display Management Switch (DMS), Target Management Switch (TMS), and Countermeasures Management Switch (CMS) --- the three HOTAS switches covered in this guide --- control display management, sensor management, and defensive systems.

SOI and Master Mode determine how HOTAS inputs are interpreted:

\begin{itemize}
    \item \textbf{Sensor of Interest (SOI):} determines which display receives HOTAS inputs.
    \item \textbf{Master Modes (NAV, A-A, A-G):} define the available weapon systems and sensor configurations.
\end{itemize}

For the complete HOTAS control reference, see \dashrefs{2.1.1.1.4 and 2.1.5}.

%============================================================================
% SECTION 2.1 — SENSOR OF INTEREST (SOI)
%============================================================================
\section{Sensor of Interest (SOI)}
\label{sec:C2-S1}

The F-16 cockpit contains three displays: the Head-Up Display (HUD) and two Multifunction Displays (MFD). The Sensor of Interest (SOI) determines which of these displays currently receives HOTAS inputs. At any moment, \textit{only one} display can be designated as SOI. The active format on the display determines how the input is interpreted.

%--------------------------------------------------------------------------
% SUBSECTION 2.1.1 — SOI FUNDAMENTALS
%--------------------------------------------------------------------------
\subsection{SOI Fundamentals}
\label{sec:C2-S1-S1}

The SOI mechanism manages \textit{where} HOTAS inputs are directed, not \textit{what} those inputs do. The specific function of a HOTAS input --- for example, whether \textbf{TMS} Up designates a target, breaks a track, or cycles through a mode --- depends on delivery modes, current Master Mode, and active sensor format.

%--------------------------------------------------------------------------
% SUBSECTION 2.1.2 — SOI SYMBOLOGY
%--------------------------------------------------------------------------
\subsection{SOI Symbology}
\label{sec:C2-S1-S2}

SOI designation is indicated by a visual marker on each display:

\begin{itemize}
	\item \textbf{HUD as SOI:} An asterisk symbol (\texttt{*}) appears in the upper left corner of the HUD, above the airspeed scale.
	\item \textbf{MFD as SOI:} A border outline is drawn around the edges of the MFD display, forming a box that distinguishes the SOI display from the non-SOI display.
	\item \textbf{NOT SOI indication:} The text \texttt{NOT SOI} appears in the center of any SOI-capable MFD format that is not currently designated as SOI. This indicator does not appear on MFD formats that cannot be designated as SOI.
\end{itemize}

%--------------------------------------------------------------------------
% SUBSECTION 2.1.3 — SOI-CAPABLE DISPLAYS
%--------------------------------------------------------------------------
\subsection{SOI-Capable Displays}
\label{sec:C2-S1-S3}

SOI designation depends on Master Mode and MFD format.

Master Mode determines SOI eligibility. MFD formats can be designated as SOI in \textit{all Master Modes}. The HUD, however, can be designated as SOI \textit{only in NAV and A-G Master Modes}.

MFD format determines SOI eligibility \textit{within} the MFD. Sensor formats that accept targeting and cursor control inputs --- such as FCR, TGP, and WPN --- can be designated as SOI. Informational and configuration formats, however --- such as DTE, TCN, and FLCS --- cannot be designated as SOI.

For SOI availability by Master Mode and MFD format, see \secref{sec:C3-S1}.

%============================================================================
% SECTION 2.2 — MASTER MODES
%============================================================================
\section{Master Modes}
\label{sec:C2-S2}

Master Mode defines the active sensor, display, and weapon configuration of the F-16 avionics suite. It also determines SOI availability and switch functions.

%--------------------------------------------------------------------------
% SUBSECTION 2.2.1 — PRIMARY AND OVERRIDE MASTER MODES
%--------------------------------------------------------------------------
\subsection{Primary and Override Master Modes}
\label{sec:C2-S2-S1}

Three primary Master Modes (NAV, A-A, and A-G) and two override modes (DGFT and MSL OVRD) are available in the F-16:

\begin{itemize}

	\item \textbf{Navigation (NAV):} The default Master Mode, active when no other mode is selected. Both air-to-air and air-to-ground sensor modes are available in NAV. Navigation symbology is displayed on the HUD.
    \item \textbf{Air-to-Air (A-A):} Selected via the A-A button on the ICP. Only air-to-air sensor modes are available in A-A. Air-to-air weapon selection and engagement symbology are displayed on the HUD.
    \item \textbf{Air-to-Ground (A-G):} Selected via the A-G button on the ICP. Both air-to-air and air-to-ground sensor modes are available in A-G. Air-to-ground weapon delivery symbology is displayed on the HUD.
    \item \textbf{Dogfight (DGFT) and Missile Override (MSL OVRD):} Selected via the Dogfight/Missile Override switch on the throttle — DGFT outboard, MSL OVRD inboard. Only air-to-air sensor modes are available in both modes. DGFT and MSL OVRD take precedence over all other Master Modes except Emergency Jettison.

\end{itemize}

HOTAS switch inputs are interpreted within the active master mode configuration --- sensor, display, and weapon options ---, which can be pre-configured via the Data Transfer Cartridge (DTC) or adjusted manually in-flight. For complete Master Mode details, see \dashref{2.1.1.2.1}.

%============================================================================
% SECTION 2.3 — DMS, TMS, AND CMS
%============================================================================
\section{DMS, TMS, and CMS}
\label{sec:C2-S3}

DMS, TMS, and CMS — including all HOTAS tables — are covered in \chapref{chap:C3}, \chapref{chap:C4}, and \chapref{chap:C5}, respectively.

%--------------------------------------------------------------------------
% SUBSECTION 2.3.1 — DISPLAY MANAGEMENT SWITCH (DMS)
%--------------------------------------------------------------------------
\subsection{Display Management Switch (DMS)}
\label{sec:C2-S3-S1}

The Display Management Switch (DMS) is a four-way, momentary hat switch on the side-stick controller. Its primary functions are SOI designation and MFD format cycling (see \dashref{2.1.5}). The DMS selects \textit{which display or sensor} receives HOTAS inputs. It does not designate targets or control sensor modes.

%--------------------------------------------------------------------------
% SUBSECTION 2.3.2 — TARGET MANAGEMENT SWITCH (TMS)
%--------------------------------------------------------------------------
\subsection{Target Management Switch (TMS)}
\label{sec:C2-S3-S2}

The Target Management Switch (TMS) is a four-way, momentary hat switch on the side-stick controller. Its primary functions are target designation, sensor management, and markpoint creation (see \dashref{2.1.5}). The function of each TMS direction depends on the Master Mode, the active sensor mode, and the current SOI display.

%--------------------------------------------------------------------------
% SUBSECTION 2.3.3 — COUNTERMEASURES MANAGEMENT SWITCH (CMS)
%--------------------------------------------------------------------------
\subsection{Countermeasures Management Switch (CMS)}
\label{sec:C2-S3-S3}

The Countermeasures Management Switch (CMS) is a four-way, momentary hat switch on the side-stick controller. Its primary functions are countermeasures dispensing and ECM control (see \dashref{2.1.5}). CMS actuation is independent of both SOI designation and Master Mode.

%============================================================================
% END OF SECTION
%============================================================================

\end{document}